% =============================================================================
% A Reproducible Pipeline for Architecture Comparison in Low-Resource
% Language Modelling
% Class final project, Natural Language Processing, Insper, 2026
% =============================================================================

\documentclass[11pt]{article}

% --- packages -----------------------------------------------------------------
\usepackage[a4paper, margin=1in]{geometry}
\usepackage[utf8]{inputenc}
\usepackage[T1]{fontenc}
\usepackage{mathptmx}                    % Times-like body + math (legacy, works everywhere)
\usepackage{microtype}                    % nicer typography
\usepackage{graphicx}
\usepackage{booktabs}
\usepackage{caption}
\usepackage{subcaption}
\usepackage[hidelinks]{hyperref}
\usepackage{xcolor}
\usepackage{titlesec}
\usepackage{titling}
\usepackage{enumitem}
\usepackage{listings}
\usepackage{authblk}
\usepackage{abstract}

% --- styling: clean, single column, Nature-ish -------------------------------
\setlength{\droptitle}{-1.2cm}
\pretitle{\begin{flushleft}\large\bfseries}
\posttitle{\par\end{flushleft}\vskip 0.4em}
\preauthor{\begin{flushleft}\normalsize}
\postauthor{\par\end{flushleft}}
\predate{\begin{flushleft}\small\itshape}
\postdate{\par\end{flushleft}\vskip 0.6em}

\titleformat*{\section}{\large\bfseries}
\titleformat*{\subsection}{\normalsize\bfseries}
\titleformat*{\subsubsection}{\normalsize\bfseries\itshape}

% Tighter spacing around section heads
\titlespacing*{\section}{0pt}{1.2em}{0.4em}
\titlespacing*{\subsection}{0pt}{0.9em}{0.3em}

% Abstract: bold heading, no italic, slightly indented body
\renewcommand{\abstractnamefont}{\normalsize\bfseries}
\renewcommand{\abstracttextfont}{\normalsize}
\setlength{\absleftindent}{0pt}
\setlength{\absrightindent}{0pt}

\setlength{\parskip}{0.4em}
\setlength{\parindent}{0pt}

% Code listings
\lstset{
  basicstyle=\ttfamily\small,
  breaklines=true,
  frame=single,
  framerule=0.4pt,
  rulecolor=\color{black!40},
  xleftmargin=0.5em,
  xrightmargin=0.5em,
  aboveskip=0.6em,
  belowskip=0.6em,
}

% =============================================================================
\title{A Reproducible Pipeline for Architecture Comparison in
       Low-Resource Language Modelling:\\
       A Case Study on xLSTM versus Transformer}

\author[1]{João~Pedro~Rodrigues~dos~Santos}
\affil[1]{Computer Engineering, Insper, São Paulo, Brazil}
\date{April 2026}

% =============================================================================
\begin{document}

\maketitle

\begin{abstract}
\noindent
Fair comparison of neural network architectures is difficult: differences in
tokenizers, hyperparameters, data preparation, and infrastructure routinely
conflate with architectural differences in published comparisons. This paper
presents an open-source pipeline designed to isolate architecture as the
sole variable in language modelling experiments at extremely small scales
(1--5\,M parameters, 1--2\,M training tokens). The pipeline standardises raw-text
ingestion, byte-level BPE tokenizer training, model construction, training,
and deterministic perplexity evaluation across user-supplied architectures.
We illustrate the pipeline by comparing xLSTM~\cite{beck2024} and a GPT-style
Transformer at two parameter scales, using nineteen 19th-century Portuguese novels
from Project Gutenberg as the training corpus. We observe a crossover:
the Transformer outperforms the xLSTM at the smallest scale, while the
xLSTM marginally surpasses the Transformer at the medium scale. We
emphasise a methodological caveat that is rarely controlled in published
comparisons: training throughput between the two architectures in our setting
differs by an order of magnitude due to implementation choices, not
architectural cost. The pipeline is released as open-source software, ready
for extension to additional architectures, corpora, and scales.
\end{abstract}

% =============================================================================
\section{Introduction}

The dominant architecture in modern language modelling is the
Transformer~\cite{vaswani2017}, which underpins both encoder
models such as BERT~\cite{devlin2019} and the GPT family of decoders.
Despite its dominance, recurrent architectures have not disappeared. The
xLSTM family, introduced recently by Beck et al.~\cite{beck2024}, revisits
the LSTM with two new components: the \emph{sLSTM}, which augments the
classical LSTM gating with exponential gates and a stabilisation mechanism;
and the \emph{mLSTM}, which replaces the scalar memory cell with a matrix
memory and is fully parallelisable. The authors argue that xLSTMs achieve
parameter efficiency competitive with Transformers in language modelling,
particularly at scale.

Most published comparisons between recurrent and attention-based language
models target the regime in which both architectures shine: large models
trained on large corpora. Considerably less is known about the
\emph{extremely low-resource regime}, where models contain a few million
parameters and are trained on a few million tokens. This regime is relevant
in two practical settings: pedagogical and prototyping experiments where
compute budgets are tight, and bootstrapping in low-resource languages
where pretraining data is fundamentally limited.

A second, methodological problem is that comparisons published across
papers are difficult to interpret because they typically differ in
tokenizer, hyperparameters, data preparation, evaluation set, and software
stack. Conclusions about architectural superiority routinely confound with
these other choices. The community has noted this reproducibility gap, but
public, end-to-end pipelines that explicitly isolate architecture as
the sole variable remain scarce.

\paragraph{Contribution.} The primary contribution of this work is
methodological: an open-source, end-to-end pipeline that isolates
architecture as the only variable in small-scale language modelling
experiments. The pipeline standardises every other component:
\begin{itemize}[leftmargin=1.2em, itemsep=0pt, topsep=0pt]
  \item raw-text ingestion and cleaning (Project Gutenberg boilerplate
        stripping);
  \item byte-level BPE tokenizer training;
  \item model factory parameterised by a single \texttt{arch} flag;
  \item training loop (AdamW, cosine schedule, gradient clipping,
        identical seeds);
  \item deterministic perplexity evaluation on a held-out corpus;
  \item logging in a single \texttt{log.jsonl} format suitable for
        cross-run aggregation and plotting.
\end{itemize}
As a secondary contribution, we apply the pipeline to xLSTM and a
GPT-style Transformer at two parameter scales (\(\sim\)1.5\,M and
\(\sim\)4.5\,M) and report empirical observations. These observations
are intentionally framed as an illustration of the pipeline rather than
as definitive claims about architectural superiority.

% =============================================================================
\section{Related Work}

\paragraph{Transformer-based language modelling.}
The Transformer~\cite{vaswani2017} is currently the default architecture
for both encoder-style models~\cite{devlin2019} and autoregressive
decoders, including the GPT line~\cite{radford2019}. The decoder
configuration we adopt as a baseline (pre-norm layers, multi-head causal
attention, GELU MLPs) is essentially the small-scale GPT recipe popularised
by Karpathy's \texttt{nanoGPT}~\cite{nanogpt}, which we use as a reference
implementation.

\paragraph{xLSTM.}
xLSTM~\cite{beck2024} comprises two block types. The sLSTM block uses
exponential gating to address the limited storage capacity of classical
LSTMs and is recurrent in nature. The mLSTM block stores associations in
a matrix, exposes a parallel formulation analogous to attention, and
provides parameter-efficient long-range memory. The original work
demonstrates competitive performance with Transformers on standard
language modelling benchmarks at scale.

\paragraph{Tokenization and tokenizer-induced confounds.}
Sennrich et al.~\cite{sennrich2016} introduced byte-pair encoding for
neural machine translation; the byte-level variant we employ is standard
in modern decoders~\cite{radford2019}. We note that comparisons across
architectures often differ in tokenizer choice, which is a known confound
since tokenizer determines effective vocabulary, sequence length, and the
intrinsic difficulty of the language modelling objective.

% =============================================================================
\section{Methodology: A Reproducible Comparison Pipeline}

The pipeline is organised in four stages
(Figure~\ref{fig:pipeline}). Stage 1 ingests raw text from a list of
Project Gutenberg book identifiers and strips the standard boilerplate
markers. Stage 2 trains a byte-level BPE tokenizer with a configurable
vocabulary size on the materialised corpus. Stage 3 constructs the model
through an architecture-agnostic factory that dispatches on a single
\texttt{arch} field in the training configuration; this is the only
configuration field that differs between comparable runs. Stage 4 runs the
training loop and produces a single \texttt{log.jsonl} per run that is
consumed downstream by analysis scripts.

\begin{figure}[h]
\centering
\fbox{\parbox{0.95\linewidth}{\centering\small
\textsf{raw books}~$\rightarrow$~\textsf{cleaning}~$\rightarrow$~\textsf{tokenizer training}~$\rightarrow$~\textsf{model factory (arch=xlstm \textbar{} transformer)}~$\rightarrow$~\textsf{training loop}~$\rightarrow$~\textsf{log.jsonl}~$\rightarrow$~\textsf{aggregated comparison}
}}
\caption{The four-stage comparison pipeline. The dispatch on
  \texttt{arch} is the only locus where architectures diverge; every
  other component is shared.}
\label{fig:pipeline}
\end{figure}

\subsection{Data}

The training corpus consists of nineteen novels from Project Gutenberg,
totalling \textasciitilde 7.5\,MB of cleaned UTF-8 text
(\textasciitilde 1.9\,M tokens at the byte-level BPE described below).
The selection covers six authors across multiple literary movements of
19th-century Brazilian and Portuguese literature: \emph{Dom Casmurro},
\emph{Quincas Borba}, \emph{Helena}, \emph{A Mão e a Luva},
\emph{Esaú e Jacó} and \emph{Memorial de Aires} by Machado de Assis
(Brazilian realism); \emph{O Primo Basílio}, \emph{O Mandarim},
\emph{Os Maias}, \emph{A Relíquia} and \emph{A Cidade e as Serras} by
Eça de Queirós (Portuguese realism); \emph{Iracema}, \emph{O Guarany}
(Vols. 1 and 2), \emph{Ubirajara} and \emph{A Pata da Gazela} by José
de Alencar (Brazilian Romantic indianism); \emph{Amor de Perdição} by
Camilo Castelo Branco (Portuguese romanticism); \emph{A Escrava Isaura}
by Bernardo Guimarães (Brazilian abolitionist romanticism); and \emph{O
Ateneu} by Raul Pompéia (Brazilian late-19th-century prose). The
held-out evaluation corpus is \emph{Memórias Póstumas de Brás Cubas}
by Machado de Assis. The holdout shares an author with six training
works (the dominant author in the period) but is itself work-disjoint;
the literature available in the public domain in Portuguese imposes
limits on the achievable holdout disjointness.

The Project Gutenberg standard boilerplate (English-language license
header and footer) is stripped using regular expressions matching the
canonical \texttt{*** START} and \texttt{*** END} markers. Without this
step, models learn to predict English legal text at the boundaries.

\subsection{Tokenizer}

We train a single byte-level BPE tokenizer with vocabulary size 4{,}000
on the same corpus that will be used for training, using the
\texttt{tokenizers} library~\cite{huggingfacetokenizers}. The byte-level
variant is chosen for two reasons: (i)~it eliminates the possibility of
\texttt{<unk>} tokens, which is desirable for accented Portuguese
orthography; and (ii)~it permits a uniform treatment across the (eventual)
multilingual extension of the pipeline. The vocabulary size of 4{,}000
is small by modern standards but appropriate for the corpus size and
deliberately stresses the language modelling objective.

\subsection{Architectures}

Two architectures are compared. Both share the same input vocabulary
size, embedding dimension, context length, and output head; they differ
only in their core block.

\paragraph{xLSTM.}
We use the official \texttt{xlstm} package~\cite{beck2024}. Each model
contains one sLSTM block and the remainder mLSTM blocks. The sLSTM block
is placed at the second position of the stack, following the recipe in
the original paper. We use the \texttt{vanilla} backend (pure PyTorch),
since the high-performance \texttt{cuda} backend requires CUDA Compute
Capability \(\ge\)\,8.0 and our experiments are conducted on CPU. This
choice is methodologically conservative for the speed comparison
discussed in Section~\ref{sec:throughput}.

\paragraph{Transformer.}
We implement a GPT-style decoder in pure PyTorch: token and learned
positional embeddings, a stack of pre-norm blocks each containing
multi-head causal self-attention (computed via
\texttt{F.scaled\_dot\_product\_attention}) and a GELU MLP, a final
layer norm, and an untied output projection. Weight initialisation
follows GPT-2~\cite{radford2019}: \(\mathcal{N}(0, 0.02^2)\) for weights and
zero for biases.

\paragraph{Parameter pairing.}
At each scale, the two architectures are configured to match in total
parameter count to within a small fraction. At the tiny scale, both
contain roughly 1.45\,M parameters (gap of 0.5\%). At the medium scale,
the xLSTM contains 4.55\,M parameters and the Transformer 4.22\,M (gap
of 7\%). The Transformer's MLP expansion factor is reduced from the
canonical~4 to~2 at the medium scale to keep the parameter budget
balanced.

\subsection{Training}

The training loop is shared across architectures and identical in every
hyperparameter except for the architecture itself. We use AdamW with
\(\beta_1{=}0.9\), \(\beta_2{=}0.95\), weight decay 0.1 (applied only
to tensors of rank $\ge$ 2, following the Chinchilla
convention~\cite{hoffmann2022}), peak learning rate \(3 \times 10^{-4}\),
linear warmup followed by cosine decay to 10\% of the peak, and gradient
clipping at norm 1.0. Sequences of length 256 are packed without padding,
using \texttt{</s>} as document separator. Batches contain 4 sequences;
no gradient accumulation is used at this scale. Training runs for 1{,}500
steps at the tiny scale and 2{,}000 steps at the medium scale. All
experiments use full \texttt{float32} precision on CPU.

\subsection{Evaluation}

Evaluation is performed at fixed intervals during training on a
deterministic, file-based holdout dataset (Section~3.1). Tokenisation,
chunking order, and the number of evaluation chunks are fixed across
runs, so reported perplexity values are directly comparable. We report
the \emph{best} dev perplexity and the \emph{final} dev perplexity, the
former because it is the value used to select checkpoints in practice,
and the latter because it characterises end-of-training behaviour.

% =============================================================================
\section{Results}

\subsection{Held-out perplexity}

Table~\ref{tab:results} reports best and final dev perplexity along with
training throughput for the four configurations. Figure~\ref{fig:scaling}
shows training loss and held-out perplexity over training steps for all
four runs.

\begin{table}[h]
\centering
\caption{Held-out perplexity on \emph{Memórias Póstumas de Brás Cubas}.
  Throughput is the mean training throughput observed during the run on
  CPU.}
\label{tab:results}
\small
\begin{tabular}{lrrrr}
\toprule
Model & Params & Best PPL & Final PPL & Throughput \\
\midrule
xLSTM tiny         & 1.46\,M & 461.9 & 514.7 & 2.8\,k tok/s \\
Transformer tiny   & 1.45\,M & \textbf{332.5} & 510.1 & 22.4\,k tok/s \\
xLSTM medium       & 4.55\,M & \textbf{229.9} & 229.9 & 1.8\,k tok/s \\
Transformer medium & 4.22\,M & 235.6 & 235.6 & 8.6\,k tok/s \\
\bottomrule
\end{tabular}
\end{table}

\begin{figure}[h]
\centering
\includegraphics[width=\linewidth]{scaling_curve.png}
\caption{Training loss (left) and held-out perplexity (right, log scale)
  for the four model--scale combinations across training steps. Tiny
  models train for 1{,}500 steps; medium models for 2{,}000 steps.}
\label{fig:scaling}
\end{figure}

Two observations stand out.

First, both architectures benefit substantially from the move from
tiny to medium scale: best held-out perplexity is reduced by approximately
50\% in both cases. This indicates that the tiny scale is undertrained
relative to the architectures' capacity, and that the held-out perplexity
is genuinely sensitive to model capacity in this regime.

Second, the relative ranking between architectures inverts with scale.
At the tiny scale, the Transformer attains a substantially lower best
perplexity than the xLSTM (332.5 vs. 461.9, a 28\% gap). At the medium
scale, the ordering inverts and the xLSTM marginally outperforms the
Transformer (229.9 vs. 235.6, a 2.5\% gap). The medium-scale gap is
small enough that, in the absence of multiple training seeds, no firm
claim of superiority can be made; we discuss this further in
Section~\ref{sec:limitations}.

A subsidiary observation is that the xLSTM medium run reached its best
held-out perplexity at the final evaluation, indicating that the run was
\emph{undertrained} at the budget chosen. The Transformer medium also
ended with best equal to final, suggesting both runs would benefit from
additional steps; a confirmatory run with the longer schedule is left
to future work.

\subsection{Throughput and the implementation caveat}\label{sec:throughput}

Table~\ref{tab:results} also reports the mean training throughput in
tokens per second. The Transformer trains roughly 5--8$\times$ faster
than the xLSTM in our setting. We emphasise that this is \emph{not} a
fundamental architectural cost. The reported numbers reflect the use of
the \texttt{vanilla} (pure PyTorch) xLSTM backend; the dedicated CUDA
kernels of the official xLSTM implementation, which parallelise the
recurrence via a custom scan operation, would close most of this gap on
GPU. Conversely, the Transformer benefits substantially in CPU from the
fused \texttt{F.scaled\_dot\_product\_attention} kernel, which exploits
SIMD/AVX in a way that the loop-based recurrent path does not.

This implementation-versus-architecture distinction is a methodological
point that the pipeline makes explicit by logging throughput in the same
\texttt{log.jsonl} as the perplexity numbers, and that we believe is
under-reported in arch-comparison papers in general.

\subsection{Qualitative inspection}

We provide samples from each trained model in Appendix~\ref{app:samples}.
Briefly, the medium-scale models produce text that is recognisably
Portuguese and reproduces the orthographic conventions of 19th-century
Brazilian literature (\emph{pae} for \emph{pai}, contractions like
\emph{n'um}, dialogue with em-dashes). The tiny-scale models produce
fragments with isolated correct words but limited syntactic coherence.
Sample-level differences between architectures are subjective and we do
not draw quantitative conclusions from them.

% =============================================================================
\section{Discussion}

The pipeline succeeds in producing comparable runs in the sense that
matters for arch comparison: every component except the architecture is
shared, every hyperparameter is logged, and the held-out evaluation is
deterministic. We believe this is a useful artefact in itself,
independently of the specific result reported. In particular, the
pipeline can be extended to (i)~additional architectures by adding a
new branch to the model factory; (ii)~larger scales by adjusting the
configuration files (a 50\,M-parameter configuration is included in the
repository as a target for future GPU experiments); and (iii)~other
corpora by replacing the book-list module.

The crossover observation in our case study is consistent with the
hypothesis advanced by Beck et al.~\cite{beck2024} that xLSTM advantages
emerge with scale. We stress, however, that the gap at the medium
scale is small (2.5\%) and that we have not measured run-to-run variance.
A definitive claim would require multiple seeds; we view our observation
as suggestive rather than conclusive.

% =============================================================================
\section{Limitations and Future Work}\label{sec:limitations}

\paragraph{Single seed.}
All numbers reported are from a single training run per configuration.
The 2.5\% gap at the medium scale is comparable in magnitude to seed
variance commonly observed in small-scale language modelling. Multiple
seeds (we suggest at least three) would be needed before drawing
quantitative conclusions about the medium-scale ranking.

\paragraph{Domain.}
The corpus consists exclusively of 19th-century Portuguese-language
literature. We do not claim our observations transfer to other domains
(modern Portuguese, technical text, code, etc.) or to other languages.
The pipeline structure makes such extensions straightforward, but they
remain to be performed.

\paragraph{Scale.}
The largest model in this study contains \(\sim\)4.5\,M parameters and
sees \(\sim\)2\,M tokens during training. This is several orders of
magnitude smaller than typical pretraining setups for which the original
xLSTM and Transformer comparisons are reported. The repository contains
a 50\,M-parameter configuration paired with a streaming Wikipedia
pipeline, intended for future GPU experiments in which the question of
architectural advantage at near-Chinchilla-optimal token budgets can
be addressed in a more rigorous fashion.

\paragraph{Holdout author overlap.}
The holdout work is by Machado de Assis, who is also the author of six
of the nineteen training works (Machado being the most prolific author
of the period in the public domain). While the works themselves are
disjoint, an author-disjoint holdout would be methodologically
preferable; it was not adopted because the held-out perplexity numbers
must remain comparable across the runs reported here, all of which used
this holdout.

\paragraph{Subtraining.}
Both medium-scale runs reached their best held-out perplexity at the
final evaluation, indicating that the chosen step budget is below the
point of diminishing returns. A longer schedule (e.g., 4{,}000 steps with
proportionally larger warmup) is the natural next step, and may shift
the medium-scale ranking.

% =============================================================================
\section{Conclusion}

We have presented an open-source pipeline that isolates architecture as
the sole variable in small-scale language modelling experiments. As an
illustrative case study, we have applied the pipeline to xLSTM and a
GPT-style Transformer at two parameter scales using a corpus of
19th-century Portuguese literature, and observed a crossover effect in
held-out perplexity: the Transformer wins at the smallest scale, and the
xLSTM marginally surpasses the Transformer at the medium scale. The
observation is consistent with the hypothesis that xLSTM advantages
emerge with scale, but the small medium-scale gap and the absence of
multiple seeds preclude strong claims. We have additionally highlighted
a methodological subtlety that is rarely controlled in published
arch-comparison work: training throughput differences between recurrent
and attention-based architectures in our setting span an order of
magnitude due to backend implementation rather than fundamental
architectural cost.

% =============================================================================
\section*{Code and Data Availability}

The full pipeline, including all configuration files, training scripts,
analysis scripts, and the four trained checkpoints from this study, is
available at: \url{https://github.com/Joao-pedrosantos/NaturalLanguageProcessing}.
All training data is publicly available from Project
Gutenberg under public-domain or non-restrictive licenses.

% =============================================================================
\begin{thebibliography}{9}

\bibitem{vaswani2017}
Vaswani, A. \emph{et al.}
Attention is all you need.
\emph{Advances in Neural Information Processing Systems 30} (2017).

\bibitem{devlin2019}
Devlin, J., Chang, M.-W., Lee, K. \& Toutanova, K.
BERT: Pre-training of deep bidirectional transformers for language
understanding.
\emph{NAACL-HLT} (2019).

\bibitem{beck2024}
Beck, M. \emph{et al.}
xLSTM: Extended long short-term memory.
\emph{Advances in Neural Information Processing Systems 37} (2024).

\bibitem{radford2019}
Radford, A. \emph{et al.}
Language models are unsupervised multitask learners.
OpenAI technical report (2019).

\bibitem{sennrich2016}
Sennrich, R., Haddow, B. \& Birch, A.
Neural machine translation of rare words with subword units.
\emph{Proceedings of ACL} (2016).

\bibitem{hoffmann2022}
Hoffmann, J. \emph{et al.}
Training compute-optimal large language models.
\emph{Advances in Neural Information Processing Systems 35} (2022).

\bibitem{nanogpt}
Karpathy, A.
nanoGPT: The simplest, fastest repository for training/finetuning
medium-sized GPTs.
GitHub repository (2022). \url{https://github.com/karpathy/nanoGPT}.

\bibitem{huggingfacetokenizers}
Hugging Face.
\texttt{tokenizers}: Fast state-of-the-art tokenizers optimised for
research and production.
GitHub repository.
\url{https://github.com/huggingface/tokenizers}.

\end{thebibliography}

% =============================================================================
\appendix
\section{Generated Samples}\label{app:samples}

All samples are generated with temperature 0.8 and top-$k$ 40, conditioned on
the prompt \texttt{"Capitu olhou para mim"}. We reproduce the model output
verbatim; line breaks reflect the model's own newline tokens.

\paragraph{xLSTM tiny (1.46\,M).}
\begin{quote}\small
\itshape
Capitu olhou para mim. E ``I?'' Eu, ``S and I said, ``All to make
--Não, com a to and we had, que me as, e a ser em lir a Cas. --V O seus
ret knowra, e os meus bo-se, com ella, a cabeça, a cim, mas, n'um que é
que o pae de Capitú ao Ca,
\end{quote}

\paragraph{Transformer tiny (1.45\,M).}
\begin{quote}\small
\itshape
Capitu olhou para mim, peu á Europa; elle me fez um de uma fóimente o
nosso velho, não teria a primeira, e Conversinaria-se á vida, que eu ia
do bureira. O APITULO No dia, de CA E a mãe, e o que ella Hou-se, quando
se os lugará
\end{quote}

\paragraph{xLSTM medium (4.55\,M).}
\begin{quote}\small
\itshape
Capitu olhou para mim a passel. Os que eu não ser não a minha mãe, a
minha mãe, em dirá á porta de dizer. Ora, mas não é preciso que eram que
não lhe e eu mesmo. Vinha-me com a verdade com seu filho, e não ha de
profeira. --Ha? Não, não ha de mim. --Aqui, mas não se não ha do outro
de Deus. --Não,
\end{quote}

\paragraph{Transformer medium (4.22\,M).}
\begin{quote}\small
\itshape
Capitu olhou para mim, peu á Europa; elle me fez um de uma fóimente o
nosso velho, não teria a primeira, e Conversinaria-se á vida, que eu
ia do bureira. O APITULO No dia, de CA E a mãe, e o que ella Hou-se,
quando se os lugará
\end{quote}

\medskip

We note qualitatively that medium-scale samples reproduce 19th-century
Portuguese orthographic conventions (\emph{pae}, \emph{á}, \emph{ella},
contractions like \emph{n'um}) and recognisable dialogue structure with
em-dashes. Tiny-scale samples produce isolated correct words but lower
syntactic coherence. We refrain from quantitative qualitative claims.

\end{document}
